% Generated from locale/tr/content/lambda-calculus/introduction/reduction.tex; do not hand-edit.


\olfileid[tr]{lam}{int}{red}
\olsection{Lambda Terimlerinin İndirgenmesi}

Lambda terimleriyle ne yapılabilir? Sadeleştirilebilirler. $M$ ve $N$
herhangi iki lambda terimi ve $x$ herhangi bir değişkense, yerine koymadan
sonra $M$'nin bağlı değişkenleri $N$'nin serbest değişkenleriyle çakışacaksa
bu bağlı değişkenleri yeniden adlandırdıktan sonra, $M$ içinde $x$ yerine
$N$ koymanın sonucunu $\Subst{M}{N}{x}$ ile gösterebiliriz. Örneğin
\[
\Subst{(\lambd[w][xxw])}{yyz}{x} = \lambd[w][(yyz)(yyz)w].
\]

\begin{digress}
Yerine koyma için alternatif gösterimler $[N/x]M$, $[x/N]M$ ve
$M[x/N]$'dir. Dikkatli olun!
\end{digress}

Sezgisel olarak $(\lambd[x][M])N$ ile $\Subst{M}{N}{x}$ aynı anlama gelir;
ilk terimi ikincisiyle değiştirme işlemine \emph{$\beta$-daraltması} denir.
$(\lambd[x][M])N$'ye \emph{indirgenebilir ifade},
$\Subst{M}{N}{x}$'e onun \emph{daraltılmışı} denir. Genel olarak bir $P$
terimi, bir alt terimin $\beta$-daraltmasıyla $P'$ terimine
dönüştürülebiliyorsa, $P$'nin \emph{bir adımda $P'$'ye $\beta$-indirgendiğini}
söyler ve $P\redone P'$ yazarız. $P$'den $P'$'ye belli sayıda (belki sıfır)
tek adımlı indirgemeyle ulaşılabiliyorsa $P$, $P'$'ye
\emph{$\beta$-indirgenir}; simgesel olarak $P\red P'$. Daha fazla
$\beta$-indirgenemeyen bir terime \emph{$\beta$-indirgenemez} ya da
\emph{$\beta$-normal} denir. Bağlam açık olduğunda ``$\beta$-indirgenir''
yerine yalnızca ``indirgenir'' vb. diyeceğiz.

Bazı örnekleri ele alalım.
\begin{enumerate}
\item Şunlar geçerlidir:
\begin{align*}
(\lambd[x][xxy]) \lambd[z][z] & \redone (\lambd[z][z])(\lambd[z][z]) y \\
& \redone (\lambd[z][z]) y \\
& \redone y.
\end{align*}
\item Bir terimi ``sadeleştirmek'' onu daha karmaşık hâle getirebilir:
\begin{align*}
(\lambd[x][xxy])(\lambd[x][xxy]) & \redone (\lambd[x][xxy])(\lambd[x][xxy])y \\
& \redone (\lambd[x][xxy])(\lambd[x][xxy])yy \\
& \redone \dots
\end{align*}
\item Bir terimi değişmeden de bırakabilir:
\[
(\lambd[x][xx])(\lambd[x][xx]) \redone (\lambd[x][xx])(\lambd[x][xx]).
\]
\item Ayrıca bazı terimler birden çok yolla indirgenebilir. Örneğin en dıştaki
  uygulamayı daraltarak
\[
(\lambd[x][(\lambd[y][yx]) z)] v \redone (\lambd[y][yv]) z,
\]
en içtekini daraltarak da
\[
(\lambd[x][(\lambd[y][yx]) z)] v \redone (\lambd[x][zx]) v
\]
elde edilir. Bununla birlikte bu durumda iki terimin de aynı $zv$ terimine
indirgenmeye devam ettiğine dikkat edin.
\end{enumerate}

Son örnekteki nihai sonuç bir rastlantı değildir; lambda kalkülüsünün
``Church--Rosser özelliği'' diye bilinen derin ve önemli bir özelliğini
gösterir.

